\subsection{Measurable Functions}

\begin{definition} \label{def:borel-and-lebesgue-measurability} \hlt{Borel versus Lebesgue Measurability}\\
The \hlt{Borel $\sigma$-algebra} on $\R^p$ is generated by open sets; its elements are \hlt{Borel sets}. A function is \hlt{Borel measurable} if inverse images of Borel sets are Borel. Lebesgue measurability contains all Borel sets and all subsets of Lebesgue null sets; this completion goes beyond the Borel $\sigma$-algebra.
\end{definition}

\begin{perspective} \label{per:borel-lebesgue-tradeoff} \hlt{Topology versus Completion}\\
Borel measurability tracks topology; Lebesgue measurability completes measure by adding subsets of null sets.\\
\hlt{Key Insight}: Borel sets are topological; Lebesgue sets are complete.
\end{perspective}

\begin{definition} \label{def:measure-space-and-measurable-functions} \hlt{Measure Spaces and Measurable Functions}\\
A \hlt{measure space} is a set $X$ together with a $\sigma$-ring $\mathcal{M}$ of subsets of $X$ and a nonnegative countably additive set function $\mu$ on $\mathcal{M}$. If $X\in\mathcal{M}$, then $X$ is a \hlt{measurable space}.\\
For $\bowtie\in\{>,\geq,<,\leq\}$, write $\{f\bowtie a\}\defeq\{x\in X:f(x)\bowtie a\}$.\\
An extended-real-valued $f$ on $X$ is \hlt{measurable} if $\{f>a\}\in\mathcal{M}$ for every $a\in\R$.
\end{definition}

\begin{proposition} \label{prop:measurable-function-open-set-test} \hlt{Open-Set Test for Euclidean Targets}\\
Let $\Omega\subseteq\R^n$ be Lebesgue measurable and let $f:\Omega\to\R^m$. The following tests are equivalent:
\begin{enumerate}[label=(\roman*)]
\setlength{\itemsep}{0pt}
\item $f^{-1}(V)$ is measurable for every open $V\subseteq\R^m$;
\item $f^{-1}(B)$ is measurable for every open box $B\subseteq\R^m$;
\item every coordinate function $f_j:\Omega\to\R$ is measurable in the threshold-set sense of Definition \ref{def:measure-space-and-measurable-functions}.
\end{enumerate}
In particular, every continuous $f:\Omega\to\R^m$ is measurable.
\begin{proof}
\begin{enumerate}[label=(\roman*).]
\setlength{\itemsep}{0pt}
\item[$(i)$]$\Rightarrow (ii)$. Open boxes are open.
\item[$(ii)$]$\Rightarrow (i)$. By Lemma \ref{lem:open-sets-countable-unions-rational-boxes}, write $V=\bigcup_kB_k$ with rational boxes $B_k$.\\
Then $f^{-1}(V)=\bigcup_kf^{-1}(B_k)$ is measurable.
\item[$(ii)$]$\Rightarrow (iii)$. By the previous implication, $f^{-1}(U)$ is measurable for every open $U\subseteq\R^m$.\\
Hence $\{f_j>a\}=f^{-1}(\{y:y_j>a\})$ is measurable.
\item[$(iii)$]$\Rightarrow (ii)$. If $B=\prod_{j=1}^{m}(a_j,b_j)$, then $f^{-1}(B)=\bigcap\limits_{j=1}^{m}\bigl(\{f_j>a_j\}\cap\{f_j<b_j\}\bigr)$,\\
where $\{f_j<b_j\}=\Omega\setminus\bigcap\limits_{\ell=1}^{\infty}\{f_j>b_j-1/\ell\}$. Thus each factor is measurable, hence so is $f^{-1}(B)$.
\end{enumerate}
If $f$ is continuous, Theorem \ref{thm:continuity-open-preimages} makes $f^{-1}(V)$ open relative to $\Omega$.\\
By Theorem \ref{thm:relative-open-intersection-with-open-set}, it is $W\cap\Omega$ for some open $W\subseteq\R^n$, hence measurable by Theorem \ref{thm:lebesgue-measure-basic-properties-checklist}.\\
\hlt{Structural Commentary}: Euclidean-valued measurability is just countable box algebra applied to inverse images.\\
\hlt{Key Insight}: Open sets reduce to rational boxes; vector functions reduce to coordinate threshold tests.
\end{proof}
\end{proposition}

\begin{theorem} \label{thm:measurable-function-closure} \hlt{Closure Properties of Measurable Functions}\\
Let $(X,\mathcal{M})$ be a measurable space.
\begin{enumerate}[label=(\roman*)]
\setlength{\itemsep}{0pt}
\item The four tests $\{f>a\},\{f\geq a\},\{f<a\},\{f\leq a\}\in\mathcal{M}$ for all $a\in\R$ are equivalent.
\item $f$ measurable implies $\abs{f}$ measurable.
\item If $f_n$ are measurable, then $\sup_nf_n,\inf_nf_n,\limsup_nf_n,\liminf_nf_n$ are measurable.
\item If $f,g$ real-valued measurable, $F:\R^2\to\R$ continuous, then $F(f,g)$ measurable; so are $f+g$ and $fg$.
\end{enumerate}
\begin{proof}
\begin{enumerate}[label=(\roman*)]
\setlength{\itemsep}{0pt}
\item Use $\{f\geq a\}=\bigcap\limits_{m\geq1}\{f>a-1/m\}$; $\{f<a\}=X\setminus\{f\geq a\}$; $\{f\leq a\}=\bigcap\limits_{m\geq1}\{f<a+1/m\}$.
\item For $a>0$, $\{\abs{f}<a\}=\{f<a\}\cap\{f>-a\}$; for $a\leq0$ the set is empty.
\item $\{\sup\limits_nf_n>a\}=\bigcup\limits_n\{f_n>a\}$. Also $\limsup\limits_nf_n=\inf\limits_m\sup\limits_{n\geq m}f_n$; $\liminf\limits_nf_n=\sup\limits_m\inf\limits_{n\geq m}f_n$; inf follow from (i).
\item For $G_a=\{(u,v):F(u,v)>a\}$, write $G_a=\bigcup\limits_k(\alpha_k,\beta_k)\times(\gamma_k,\delta_k)$.\\
Then $\{F(f,g)>a\}=\bigcup\limits_k\bigl(\{\alpha_k<f<\beta_k\}\cap\{\gamma_k<g<\delta_k\}\bigr)\in\mathcal{M}$.
\end{enumerate}
\hlt{Structural Commentary}: Measurability is countable set algebra applied to level sets.\\
\hlt{Key Insight}: Once threshold sets are measurable, ordinary limiting and algebraic operations stay measurable.
\end{proof}
\end{theorem}

\begin{definition} \label{def:simple-functions-measurable} \hlt{Simple Functions}\\
Let $(X,\mathcal{M})$ be a measurable space. For $E\subseteq X$, set $\mathbf{1}_E(x)\defeq
\begin{cases}
1,&x\in E,\\
0,&x\notin E.
\end{cases}$. Real-valued $s$ is \hlt{simple} if its range is finite.\\
With distinct values $c_i$ and sets $E_i=\{s=c_i\}$, $s=\sum\limits_{i=1}^{n}c_i\mathbf{1}_{E_i}$ and $s$ is measurable iff $E_i\in\mathcal{M}$ for all $i$.
\end{definition}

\begin{theorem} \label{thm:simple-approximation-measurable-functions} \hlt{Approximation by Simple Functions}\\
Every real $f$ on $X$ is the pointwise limit of simple functions $s_n$.\\
If $f$ is measurable, choose $s_n$ measurable; if $f\geq0$, choose $0\leq s_n\uparrow f$.
\begin{proof}
For $f\geq0$, set $E_{ni}=\left\{x:\frac{i-1}{2^n}\leq f(x)<\frac{i}{2^n}\right\}$; $F_n=\{f\geq n\}$; and $s_n=\sum\limits_{i=1}^{n2^n}\frac{i-1}{2^n}\mathbf{1}_{E_{ni}}+n\mathbf{1}_{F_n}$.\\
Then $s_n$ is simple, measurable when $f$ is, and $0\leq s_n\uparrow f$. For general $f$, approximate $f^+$ and $f^-$ and subtract.\\
\hlt{Structural Commentary}: Simple functions quantize the range of $f$.\\
\hlt{Key Insight}: Dyadic range partitions give pointwise approximation while preserving measurability.
\end{proof}
\end{theorem}

\begin{definition} \label{def:almost-everywhere-and-null-set-invariance} \hlt{Almost Everywhere}\\
For measurable $E$, write $f\sim_{\mathrm{a.e.}}g$ on $E$ if $\mu(\{x\in E:f(x)\neq g(x)\})=0$. A property holds \hlt{almost everywhere} (a.e.) on $E$ if its failure set in $E$ has measure zero. Then $\sim_{\mathrm{a.e.}}$ is an equivalence relation; if $f\sim_{\mathrm{a.e.}}g$ on $E$, integrals over measurable $A\subseteq E$ agree whenever they exist. If $f\in\mathscr{L}^1(\mu)$ on $E$, then $f$ is finite a.e. on $E$.
\end{definition}

\begin{definition} \label{def:almost-everywhere-and-in-measure-convergence} \hlt{Convergence Almost Everywhere and in Measure}\\
A sequence $f_n$ \hlt{converges almost everywhere} to $f$ on $E$ if $f_n(x)\to f(x)$ for all $x\in E\setminus N$ for some null set $N$.\\
On a measure space, $f_n$ \hlt{converges in measure} to $f$ on $E$ if for every $\varepsilon>0$, $\mu\{x\in E:\abs{f_n(x)-f(x)}>\varepsilon\}\to0$.
\end{definition}

\begin{method} \label{meth:measure-measurability-proof-patterns-source-todo} \hlt{Measure Proof Patterns}
\begin{enumerate}[label=(\roman*)]
\setlength{\itemsep}{0pt}
\item \hlt{Disjointize}. If $A_1\subseteq A_2\subseteq\cdots$, set $B_1=A_1$ and $B_n=A_n\setminus A_{n-1}$.\\
Then $A_n=\bigcup\limits_{k=1}^nB_k$ and $\bigcup\limits_nA_n=\bigcup\limits_nB_n$ are disjoint decompositions.
\item \hlt{Approximate}. Prove identities on $\mathcal{E}$; if $A_n\to A$ and $B_n\to B$ in $d_\mu(A,B)=\mu^\ast(A\triangle B)$, pass through the symmetric-difference estimates for $\cup,\cap,\setminus$.
\item \hlt{Threshold}. To prove real-valued $f$ measurable, verify $\{f>a\}\in\mathcal{M}$ for all $a\in\R$; for Euclidean-valued maps, use Proposition \ref{prop:measurable-function-open-set-test}.
\item \hlt{Topological bridge}. Use Theorem \ref{thm:continuity-open-preimages}: open inverse images reduce continuous constructions to countable unions of interval or rectangle preimages.
\end{enumerate}
\hlt{Key Insight}: Measure theory alternates between countable set algebra and approximation by elementary sets.
\end{method}
